\documentclass[11pt,a4paper]{article}

\usepackage[margin=2cm]{geometry}
\usepackage[T1]{fontenc}
\usepackage{lmodern}
\usepackage{graphicx}
\usepackage{booktabs}
\usepackage{tabularx}
\usepackage{amsmath}
\usepackage{caption}
\usepackage{enumitem}
\usepackage{listings}
\usepackage{xcolor}
\usepackage[hidelinks]{hyperref}

\captionsetup{font=normalsize, labelfont=bf}
\setlength{\parskip}{0.4em}
\setlength{\parindent}{0pt}
\setlist{itemsep=0.15em, topsep=0.3em}

\lstset{basicstyle=\ttfamily\small, breaklines=true, frame=single,
        commentstyle=\color{gray}, columns=fullflexible}


\begin{document}
\begin{center}
  {\Large\bfseries Midland Premier League 2023/24: A Statistical Analysis for Ashford United FC\par}
  \vspace{0.4em}
  Group M \quad|\quad [Group names] \quad|\quad Nov. 2027
\end{center}



\section{Introduction}

Ashford United FC has commissioned Meridian Analytics to analyse data from the
2023/24 Midland Premier League season. The data are simulated. They consist of
records for 620 players across the league's 24 clubs and the results of Ashford
United's 46 league matches. Our aim is to identify and quantify relationships in
these data that are useful to the club. We address three questions:
\begin{enumerate}[label=\textbf{Q\arabic*.}, leftmargin=2.6em]
  \item \textbf{...}

  \item \textbf{...} 
  
  \item \textbf{...}
  
\end{enumerate}



\section{Data and exploratory analysis}\label{sec:data}

Table~\ref{tab:vars} lists the variables in the two datasets. Neither dataset
contains missing values or duplicated records.

\begin{table}[ht]
\centering
\caption{Variables in the players dataset (left; one row per player, 24--28
players per club) and the matches dataset (right; one row per Ashford United
match). Ranges are given in brackets.}
\label{tab:vars}
\begin{minipage}[t]{0.55\textwidth}
\vspace{0pt}
\begin{tabularx}{\linewidth}{@{}l X@{}}
\toprule
\multicolumn{2}{@{}l}{\textbf{Players} ($n = 620$)}\\
\midrule
\texttt{team}            & Club (24 clubs) \\
\texttt{position}        & GK, DEF, MID or FWD \\
\texttt{age}             & Age in years (17--36) \\
\texttt{minutes\_played} & Minutes played (105--3960) \\
\texttt{training\_load}  & Weekly intensity, 1 (low) to 5 \\
\texttt{prev\_injuries}  & Injuries in previous 3 seasons (0--4) \\
\texttt{bmi}             & Body mass index (19.0--27.9) \\
\texttt{injured}         & Injured in season: 1 yes, 0 no \\
\bottomrule
\end{tabularx}
\end{minipage}\hfill
\begin{minipage}[t]{0.42\textwidth}
\vspace{0pt}
\begin{tabularx}{\linewidth}{@{}l X@{}}
\toprule
\multicolumn{2}{@{}l}{\textbf{Matches} ($n = 46$)}\\
\midrule
\texttt{home\_away}      & Home or Away \\
\texttt{goals\_scored}   & By Ashford (0--8) \\
\texttt{goals\_conceded} & By Ashford (0--3) \\
\texttt{opponent\_rank}  & Final league position of opponent (1--24) \\
\bottomrule
\end{tabularx}
\end{minipage}
\end{table}

Ashford United played each of the other 23 clubs once at home and once away.
Rank~9 is the only rank missing from the opponents, so Ashford finished ninth.
Because every opponent appears once at each venue, the effects of venue and
opponent strength can be estimated separately. In the players data, 16 players
(12 of them goalkeepers) recorded the maximum of 3960 minutes, so
\texttt{minutes\_played} is capped.

\subsection{Ashford United's matches}
Ashford United's results differed markedly by venue. At home they won 15, drew
5 and lost 3 of 23 matches. Away they won 7, drew 6 and lost 10. A chi-squared
test of independence gives $X^2 = 6.77$ on 2 degrees of freedom ($p = 0.034$),
so at the 5\% level we reject the hypothesis that result and venue are
unrelated. Ashford averaged 1.74 goals per match at home and 0.87 away. They
conceded 0.57 per match at home and 1.09 away.

Figure~\ref{fig:goals} shows that Ashford tended to score more and concede fewer
goals against clubs lower in the table. The Spearman correlations with opponent
rank are $0.29$ for goals scored and $-0.32$ for goals conceded. These are
modest associations, since goals in a single match are highly variable. The
fitted curves for the two venues have similar slopes, which suggests that the
association with opponent strength does not depend on venue. We test this in
Section~4.

\begin{figure}[ht]
\centering
\includegraphics[width=\textwidth]{fig1_goals_vs_rank.pdf}
\caption{Goals scored (left) and conceded (right) by Ashford United in each
match, against the opponent's final league rank. Points are separated slightly
in the vertical direction to show repeated values. The lines are Poisson
regression curves fitted separately for each venue, shown as a visual guide.}
\label{fig:goals}
\end{figure}

Both goal counts are small whole numbers, mostly 0 to 2, which suggests a
Poisson model, in which the variance equals the mean. For goals conceded the
variance (0.64) is slightly below the mean (0.83). For goals scored the
variance (1.99) is well above the mean (1.30). This is almost entirely due to
one match, an 8--0 home win against the bottom club. Without it, the variance
(1.00) and mean (1.16) are close. A preliminary Poisson fit shows that this
match has a large influence on the estimated effects, so Section~4 reports the
analysis with and without it.

\subsection{Player injuries}
Of the 620 players, 115 (18.5\%) were injured during the season.
Figure~\ref{fig:injury} shows how the proportion injured varies with four
characteristics. It rises steadily with the number of previous injuries, from
11\% of players with none to 43\% of those with four. It is about 26\% for
training loads of 4 or 5, against 11--16\% for loads of 1 to 3. It is 34\% for
players over 29, against about 14\% for those aged 25 or under. Chi-squared
tests for a linear trend in proportions give strong evidence of an increasing
trend for each of these variables ($X^2 = 35.1$, $11.4$ and $13.3$ on 1 degree
of freedom; all $p < 0.001$). Goalkeepers have the highest proportion (27\%),
but there are only 62 of them, and a chi-squared test gives no evidence of an
association between position and injury ($X^2 = 4.38$ on 3 degrees of freedom,
$p = 0.22$). Injured and uninjured players had similar minutes played and BMI.

\begin{figure}[ht]
\centering
\includegraphics[width=\textwidth]{fig2_injury_rates.pdf}
\caption{Proportion of players injured during the season by previous injuries,
training load, age group and position, with 95\% confidence intervals.
%%%
The number above each interval is the group size. The dashed line is the overall
proportion (18.5\%).}
\label{fig:injury}
\end{figure}

The explanatory variables are close to uncorrelated: the largest correlation is
0.21, between training load and minutes played. Collinearity is therefore not
a concern when these variables are modelled together. Injury proportions vary
across clubs from 8\% to 44\%, with Ashford at 16\%. A chi-squared test gives
little evidence of a difference between clubs ($X^2 = 25.8$ on 23 degrees of
freedom). Players at the same club may still
share unmeasured conditions, such as facilities and medical staff, so in Q1 we
check whether a club effect changes the conclusions.


\section{Q1:}


\section{Q2:}


\section{Q3:}


\section{Conclusions and limitations}



\clearpage
\appendix
\section*{Appendix: R code}

Shared workspace: \url{[link to shared workspace]}

% Export the R code from the Rmd with knitr::purl("EDA.Rmd") (and the analysis
% files) and include it here, e.g.
% \lstinputlisting[language=R]{EDA.R}
\textit{[Complete R code to be inserted.]}

\end{document}
